\minitoc 

% ========================================================================================================
% Introduction 

L'algorithme de Metropolis-Hasting est une procédure permettant de transformer une chaîne de Markov par le rejet 
sélectof de certains de ses pas en une nouvelle chaîne admettant une mesure invariante fixée. 

\begin{proposition}[Algorithme] 
    Soit $(X_n)_{n \in \N}$ une chaîne de Markov définie à partir d'une application 
    $f : E \times [0,1] \longrightarrow E$ telle que pour tout $n \in \N$, on ait : 
        $$ X_0 \in E, \quad \text{et} \quad X_{n+1} = f(X_n, U_{n+1}) $$ 
    où $(U_n)_{n \geqslant 1}$ est une suite de var idd de loi $ \mathcal{U}([0,1]) $. 
    
    Soit $(V_n)_{n \in \N}$ une suite de var iid indépendante de $(X_n)_{n \in \N}$ et de loi $ \mathcal{U}([0,1])$. 
    Soit $ \pi$ une mesure invariante cible. On définit la chaîne de Markov $(Y_n)_{n \in \N}$ par : 
        $$ Y_{n+1} = 
            \begin{cases} 
                f(Y_n, U_{n+1}) \quad \text{si} \quad V_n < \frac{ \pi(y) Q(y,x)}{ \pi(x)Q(x,y)} \\ 
                Y_n \quad \text{sinon} 
            \end{cases} $$ 
    où $x = Y_n$ désigne l'état courant de la nouvelle chaîne au temps $n$ et $y = f(Y_n, U_{n+1})$ 
    est le nouvel état candidat au temps $n+1$ et : 
        $$ \forall x \in E, \quad Q(x, A) = \myP(f(x,U) \in A \mid Y_n = x) $$ 
\end{proposition}


% ========================================================================================================
% Cas d'un espace d'états fini 

\section{Cas d'un espace d'états fini}

\begin{proposition}[Matrice de transition]
    Soit $(X_n)_{n \in \N}$ une chaîne de Markov sur un espace d'états fini $E$. On note : 
        $$ \forall x,y \in E, \quad Q(x,y) = \myP(X_1 = y \mid X_0 = x) $$
    Supposons que $(X_n)_{n \in \N}$ est irréductible. Alors l'algorithme de Metropolis-Hasting 
    transforme la chaîne $(X_n)_{n \in \N}$ en une chaîne $(Y_n)_{n \in \N}$ de matrice de transition : 
        $$ P(x,y) = 
            \begin{cases}
                Q(x,y) \times \min \left( 1, \frac{ \pi(y) Q(y,x)}{ \pi(x) Q(x,y)} \right) \quad \text{si} \quad n \not = y \\ 
                1 - \sum_{z \not = x} P(x,z) \quad \text{si} \quad x = y  
            \end{cases}
        $$ 
\end{proposition}

\begin{proof}
    Soient $x,y \in E$ tels que $x \not = y$. Calculons $ \myP(Y_1 = y \mid Y_0 = x)$. On a : 
        \begin{align*}
            \myP( Y_1 = y \mid Y_0 = x) &= \myP_x \left( f(x,U_1) = y, V_1 < \frac{ \pi(y) Q(y,x)}{ \pi(x) Q(x,y)} \right)  \\
            &= \myP_x(X_1 = y) \myP \left( V_1 < \frac{ \pi(y) Q(y,x)}{ \pi(x) Q(x,y)}\right) \\
            &= Q(x,y) \min \left( 1, \frac{ \pi(y) Q(y,x)}{ \pi(x) Q(x,y)} \right)
        \end{align*} 
    De même : 
        \begin{align*}
            \myP(Y_1 = \mid Y_0 = x) &= 1 - \sum_{y \not = x} \myP(Y_1 = y \mid Y_0 = x) \\
            &= 1 - \sum_{y \not = x} P(x,y)
        \end{align*}
\end{proof}

\begin{proposition}
    Si pour tout $x,y \in E$, on a : 
        $$ Q(x,y) > 0 \iff Q(y,x) > 0 $$ 
    alors la mesure $ \pi$ est réversible pour la chaîne $(Y_n)_{n \in \N}$. 
\end{proposition}

\begin{proof}
    Soient $X,y \in E$. On souhaite vérifier que : 
        $$ \pi(x)P(x,y) = \pi(y)P(y,x) $$ 
    Distinguons deux cas : 
    \begin{itemize}
        \item \underline{si $x = y$} alors on a évidement $\pi(x)P(x,x) = \pi(x)P(x,x)$. 
        \item \underline{si $x \not = y$} alors : 
            $$ \pi(x) P(x,y) = \pi(x) Q(x,y) \min \left(1, \frac{ \pi(y) Q(y,x)}{ \pi(x) Q(x,y)}  \right) $$ 
        de même : 
            $$ \pi(y) P(y,x) = \pi(y) Q(y,x) \min \left(1, \frac{ \pi(x) Q(x,y)}{ \pi(y) Q(y,x)}  \right) $$ 
        Supposons que $  \pi(x) Q(x,y) \leqslant \pi(y) Q(y,x)$ alors : 
            $$ \pi(x)P(x,y) = \pi(x) Q(x,y) $$ 
        et 
            \begin{align*}
                \pi(y) P(y,x) &= \pi(y) Q(y,x) \frac{ \pi(x) Q(x,y)}{ \pi(y) Q(y,x)} \\ 
                &= \pi(x) P(x,y) 
            \end{align*}
        Si $ \pi(x) Q(x,y) > \pi(y) Q(y,x) $, le raisonnement est le même.
    \end{itemize}
\end{proof}

\begin{example}
    Soit $(X_n)_{n \in \N}$ la marche aléatoire simple sur $\Z$. Soit $ \pi$ une mesure de probabilité sur $ \Z$. 
    On cherche à construire une chaîne de Markov $(Y_n)_{n \in \N}$ sur $\Z$ avec l'algorithme de Metropolis-Hasting 
    de mesure réversible $ \pi$. On a donc les algorithmes suivants : 

    \begin{center}
        \begin{minipage}{0.45\textwidth}
            \begin{lstlisting}[language=Python]
def etapeMa(x): 
    if rand() < 1/2: 
        return x+1
    return x-1
            \end{lstlisting}
        \end{minipage}
        \hfill 
        \begin{minipage}{0.45\textwidth}
            \begin{lstlisting}[language=Python]
def algoMetHast(x): 
    y = etapeMa(x)
    if rand() < pi(y) * (1/2) / pi(x) * (1/2): 
        return y 
    return x
            \end{lstlisting}
        \end{minipage}
    \end{center}

    On a donc : 
        $$ Q(x,x+1) = Q(x,x-1) = \frac{1}{2} $$ 
        $$ \forall x,y \in E, \; |x-y| \not = 1, \quad Q(x,y) = 0 $$ 
    Et : 
        $$ P(x,x+1) = \frac{1}{2} \min \left( 1, \frac{ \pi(x+1)}{ \pi(x)}  \right) \quad \text{et} \quad P(x, x-1) = \frac{1}{2} \min \left(1, \frac{ \pi(x-1)}{ \pi(x)} \right) $$
\end{example}

\begin{prop}
    Si $Q$ est irréductible et 
        $$ \forallQ(x,y) > 0 \iff Q(y,x) > 0 $$ 
    et $ \pi(x) > 0$ 
\end{prop}